\chapter[Sermon 30. Two Seals Are Opened, and the Direct Course of God's Word Is Set Forth, and a Cruel Course of Wars Against the Disobedient.]{Two Seals Are Opened, and the Direct Course of God's Word Is Set Forth, and a Cruel Course of Wars Against the Disobedient.}
\label{sermon:030}
\markright{Sermon 30. Two Seals Are Opened, and the Direct Course of God's ...}

And I saw when the Lamb opened one of the seals, and I heard one of the four beasts say, as it were the voice of thunder: Come and see. And I saw, and behold a white horse: And he that sat on him had a bow, and a crown was given unto him: And he went forth conquering, and to overcome. And when he opened the second seal, I heard the second beast say: Come and see. And there went out another horse that was red, and power was given to him that sat there on, to take peace from the earth, and that they should kill one another. And there was given unto him a great sword.

Up to this point, the Apostle has prepared the audience to hear with a calm mind about God’s judgments and the church’s destined fate, and to patiently endure all adversity, and to worship him in all things, and to give glory to his name. Consequently, he explains in a very orderly fashion God’s judgments and the church’s destiny, showing how the Son of God governs God’s ordinances and his eternal providence. And this is, as it were, a prophecy for all times and ages until the end of the world. For we should not think that only the events of one age or two are recounted here, but of all. First, all things are generally described by parts; afterward, particularly, when we come to the opening of the seventh seal. The sum is, the Lord sends forth the preaching of the truth into the world, which when people refuse and despise, they are destroyed with wars and other innumerable calamities.

\index{John, Saint}But above all things, Saint John is stirred and eager (and in him, all of us) to be attentive. And one, that is to say, the first of the beasts, excites him. One of the Sabbath is set for the first day of the week, which is truly Sunday. The voice of the beast is like unto thunder, signifying that great and weighty matters are being treated. For most great and terrible things follow, which shake the whole world. Therefore let us not play the sleepy sluggards, let us not be blind and deaf. Doubtless the slothfulness of our time is such that we little consider the works of God and what is done in our time. The storks, swallows, and doves, and the rest of living things pass us by, which carefully observe their time. Therefore, let us be stirred up here, that we should not be slothful, but should mark what things are declared and shown to us by the Lord.

And when John had carefully recorded what was done, he sees the Lamb, Christ, our Redeemer, open one seal, that is to say, the first. Immediately a white horse came forth, on whom he who sat had a bow bent, and an arrow in it. To him was given a crown, and he went forth conquering, that he might overcome. This is the vision: the explanation of which is easy. For the Lord says that he will declare the destinies of the church.

\index{Zechariah (Book of)}Horses of sundry colors are also brought forth by Zechariah in the first chapter. They signify the variable course and state of the people of Israel. The white color is consecrated to innocence, purity, victory, and felicity. Therefore, the white horse signifies the fortunate proclamation of God’s word, or the prosperous preaching of the Gospel. For upon the horse sits a horseman, who guides the horse and has a bow. Certainly, Christ promotes the course of preaching the Gospel. And Psalm 45 attributes to the same shafts or arrows, for he strikes his enemies far off and brings them into his subjection. Briefly, with the word of his mouth, he subdued to himself people and nations. Isaiah, in chapter 49, bringing in Christ speaking, says: “And he put my mouth as a sharp sword, the shadow of his hand covered me, and he put me as a sharpened arrow, he hid me in his quiver.” Through Christ, therefore, proceeds the preaching of the word; he gives strength to the preaching; he draws back his bent bow. Whatever force the word has, that is wholly due to the Horseman.

A crown is given to him, namely, a kingdom and all power of ruling. For David prophesying beforehand said, “The Lord will send forth the rod of his power out of Zion, to rule among your enemies.” Moreover, a crown is given to him, that he may crown those who serve him faithfully. And it is a figure of speech, and he went forth conquering, that he might overcome: for he who went forth is a conqueror, and for this purpose went forth, that he might overcome. For it signifies that Christ will advance the preaching of his word throughout the world, no one being able to resist, and even in defiance of the gates of Hell. For the word of the Lord endures forever.

\index{Papacy}\index{Rome}And this place teaches that the Church shall always be in the world, and likewise always the truth preached, even though the enemies’ hearts burst. But if we read over the history of the Church, we shall better understand all things, and shall perceive that this prophecy has always been most certain. Christ was once through the ministry of the word showed to the world by the Apostles, and the matter proceeded most favorably, however much the most powerful people of this world resisted it. The thing is wonderful, if one considers those five hundred years which immediately after the incarnation of our Lord are accounted. In them went forth the conqueror that he might overcome: and overcame in deed, the whole world receiving Christ and worshipping him. Since those years, as before also, certain seeds of errors began to be sown abroad. The bishops began to contend for supremacy, and who should be the universal head of the Church on earth: they began to reason about the use of images in the Church, and brought them into churches in deed, as also they called the bishop of Rome the supreme and general head of the Church on earth. And mighty princes, and in a manner the whole state of learned men conspired in these opinions: but he has vanquished, which went forth that he might vanquish. He had in his church innumerable, who bowed not their knees before this Baal. A thousand years after the incarnation of Christ, the bishops began to profanely pollute the Lord’s Supper and other undefiled doctrines of faith: but what, I pray you, did they prevail by so many councils, determinations, and earnest endeavors? He that went forth to overcome, has overcome. That white Horse has stoutly advanced to the salvation of many. For how great battles the godly and learned have sustained against the Popes and bishops in these last five hundred years, history bears witness. At this day also appears throughout the whole world, how favorably yet that white horse goes forward, which has pierced even until our time. The Gospel is believed, and that faith cannot be extinguished with any waters or fires.

\index{Antichrist}\index{Luther, Martin}You make exception that they were heretics who resisted the bishop and See of Rome in these last 500 years, such as Bertram, John Scot, nicknamed Duns, Berengarius, Arnoldus Brixianus, Waldo, Wycliffe, and Hus, Luther, and Zwinglius, and others of the same sort; moreover, some of these were also overcome and put to death by the Pope. I answer that, as men, they might err in many things; but in those things wherein with the Scripture they consent against the See of Rome, I affirm that they erred not, but spoke the truth. Whereupon it is certain that Christ overcame by them. When Micah, Elijah, Zechariah, Amos, Jeremiah, and others preached by the word of God against idols and worshippers of idols, they also were condemned as seditious and heretics; indeed, some of them were removed from the way. But was the truth vanquished? Antichrist is said to have good fortune, and to punish and afflict the strong and the people of God; but men being ministers may be oppressed, the ministry never decays. Saint Paul says that he is bound for the Gospel’s sake, but the Gospel is not to be bound. Therefore, he has overcome hitherto, and shall overcome still, which went forth that he might conquer. They stumble upon this conqueror as at the stone of offense, whoever and whatever they may be, who seek to interrupt the plain course of the Gospel.

\index{Ezekiel (Book of)}Moreover, concerning the time when the second seal should be opened, the second beast exhorts John again to attentiveness, and that we should consider what is presented to us. And now comes forth the red horse, whose color is somewhat like fire; there sits also on him a rider, to whom power is given to disturb peace on earth, and that men should kill one another. For there is given him a great sword. The red horse signifies the state of wars, full of fire and blood. He who sits on this horse is Mars, or rather the father of lies, that is, the Devil, who was a murderer from the beginning. He gathers to him the dregs of men to make civil commotions, for the wars, destruction, burning, slaughter, and desolation. You see from whence the breaking of peace is, which God hates. And we hear how it is given him: Mark given, by the just judgment of God to be permitted, that troubling all peace, he should take it away, and set men together by the ears, that one may wound and kill another. For so we read in the first book of Job, how Satan had power given him by God against Job. To bloody soldiers is given a great sword, great power to hurt, a wonderful force of fighting, as also Nahum explains it. Neither is it a rare thing in the scriptures for monarchs, tyrants, and mighty men of war to be called a sword. For so Ezekiel called Nebuchadnezzar; and Isaiah called Sennacherib, king of Assyria, a whetstone.

And the chiefest righteousness is to give every man his own. Therefore, this passage justly ascribes that which is good to God and that which is evil to the Devil. But you say, if God permits, the same that he does not prohibit he does. He does not prohibit war, for because justice will not suffer him so to do; but he commands him by war to punish the wicked and to try the good. But in permitting wars, God offends nothing, saying that for most just causes he permits the same. For they would not embrace peace offered them by the preachers of the Gospel, therefore were they worthy to be entangled with wars. The Jews did not know the day of Christ’s visitation, therefore they were justly visited by the Romans and destroyed. And this thing is in the world perpetually, that they who will not obey the Gospel must obey the Captain of the wars; they who will not hear Christ must hear Antichrist. You may not contend with God, which he does this and permits that. Worship God rather, as you have been taught in the fourth and fifth chapters.

\index{Turks}Let us examine stories and see if such wars are found, wherein men have slain themselves with mutual wounds, and have killed one another like beasts. If you will read Herodotus, Orosius, and other good historians, you may find that the Roman Emperors were troubled with most grievous wars, for no other cause than that they refused peace offered to them by the gospel. For no other cause was Rome itself at the last taken by the West Goths, and burnt and destroyed by the East Goths. The Lord had given them Christian princes, but they loved idols more. For Simmachus, governor of the city, was so bold as to demand the restoration of idolatry. I speak nothing now of Attila, nothing of the Persian and African wars. And when there was a great strife among the bishops about supremacy, the Saracens arose and became mighty. After the thousand years began the holy war, which was as bloody as it was long-lasting. Never was such a war made in all the world. Boniface VIII instituted first the year of Jubilee, a most wicked man, who also exhibited himself to be seen by the people, both as Pope and Emperor. But in the same year of one thousand three hundred, Ottoman arose in Asia, the scourge of God, the founder of the emperors of the Turks who reign at this day. For so when Solomon built places of idolatry, his enemies arose, who wonderfully vexed and afflicted the kingdom of Solomon. What wars are made nowadays, and what are the causes of wars, all wise men see. We will not receive the peaceable gospel; it is reasonable therefore that Turkish armies should invade us, that we may both feel Antichrist to be a stout warrior, and may all abhor and detest him.

But what else remains here, than that, being converted to God through Christ, we may serve the Lord in sincere faith and holy purity? For unless we repent, the groundwork is laid at the tree’s root, \&c.
